\section{From pool span to rendered frame}
\label{sec:render}

\subsection{Stable draw entries}

The renderer is a \emph{slotted} indirect renderer.  Each chunk owns one
stable draw-entry index; the entry is never compacted or renumbered while
the chunk is alive.  The packed pools are separate:
\begin{itemize}
\item a single large vertex pool (device-local, 64-byte elements);
\item a single large index pool (device-local, 32-bit elements);
\item a host-visible indirect-command array;
\item a host-visible bounds/LoD-meta array.
\end{itemize}
A chunk's vertex and index data live in a variable-size span allocated by
\code{PackedSpaceAllocator}.  The draw entry stores the span's base vertex,
first index, index count, and instance count.

\begin{figure}[htbp]
\centering
\begin{tikzpicture}[font=\small]
  \fill[accent!10] (0,0) rectangle (10,0.9);
  \foreach \x/\w/\lab in {0/2.2/chunk A,2.2/1.4/chunk B,3.6/2.8/chunk C,6.4/1.1/free,7.5/2.5/chunk D}
    \fill[accent!28] (\x,0) rectangle (\x+\w,0.9)
      node[midway,font=\scriptsize]{\lab};
  \draw (0,0) rectangle (10,0.9);
  \node[above] at (5,0.95) {vertex/index pool (one large allocation)};
  % draw entries
  \foreach \x/\lab in {0/A,2.2/B,3.6/C,7.5/D} {
    \draw (0.1+\x*0.8,1.4) rectangle (0.1+\x*0.8+0.45,1.9);
    \node[font=\scriptsize] at (0.1+\x*0.8+0.22,1.65) {\lab};
  }
  \draw[-{Latex[length=2mm]}] (0.3,1.4)--(0.5,0.95);
  \draw[-{Latex[length=2mm]}] (2.0,1.4)--(2.6,0.95);
  \draw[-{Latex[length=2mm]}] (3.2,1.4)--(4.6,0.95);
  \draw[-{Latex[length=2mm]}] (6.2,1.4)--(8.2,0.95);
  \node[below] at (5,-0.15) {stable draw entries (one per chunk slot)};
\end{tikzpicture}
\caption{Packed pools and stable draw entries.  Chunk spans are packed
best-fit into one large vertex pool and one large index pool; the draw
entry is a stable index into the indirect-command array.  Rebuilding a
chunk replaces only its span.}
\label{fig:pool}
\end{figure}

\subsection{The draw command}

A draw command is
\begin{equation}
\begin{aligned}
  \text{DrawIndexedIndirectCount}(&\text{indexCount},\text{instanceCount}=1,\\
  &\text{firstIndex},\text{vertexOffset}=\text{baseVertex},\\
  &\text{firstInstance}=\text{slot}).
\end{aligned}
\end{equation}
The vertex shader uses \code{gl_InstanceIndex} only for the per-chunk
vegetation/impostor path; the solid terrain vertex is a direct vertex
attribute.  The index stream is chunk-local: every index is a slot number
in the chunk's vertex span, and \code{vertexOffset} rebases it at draw time.

\subsection{Depth and offscreen composition}

Vegetation is rendered to its own color and depth targets on a parallel
async command buffer; the solid pass has a separate depth target.  The
composite pass tests the vegetation depth against the solid scene depth.
This decoupling is why the vegetation pass can use a different render
scale and why it can be skipped on steady-empty frames.  The solid chunks
themselves are drawn in the main solid pass, with a depth prepass and a
shadow cascade.

\subsection{Shadows}

The shadow renderer reuses the same slots and the same hierarchy.  A
cascade cull emits one compacted command per chunk per cascade, reusing the
already-selected LoD stamp from the main cull.  The shadow vertex shader
reads the same 64-byte vertex and index pool.  No second mesh exists.
The cascade copy lists and descriptor-buffer mirror include the brush map
and wind field so the shadow tessellation stage can displace and classify
the surface without a separate descriptor.

\subsection{Ray tracing}

The hybrid ray tracer builds a BLAS over the merged vertex/index pools.  A
chunk's geometry span is referenced directly:
\begin{itemize}
\item vertex buffer = the pool, vertex stride 64;
\item index buffer = the pool, first index = span's first index;
\item geometry flags = opaque or alpha-tested, depending on the material.
\end{itemize}
Because the span is published only after the finalize fence, the BLAS build
never sees partially written vertices.  A chunk rebuild requests a BLAS
refresh; the proxy boxes are used for the water reflection path and for
the fallback when the BLAS is not ready.

\subsection{Material and mapping}

The vertex carries brush, HSV, texture type, and mapping data.  The solid
fragment shader selects the material texture from the brush and the
triplanar/projection mapping from the texture type.  Because the CPU never
computed a vertex, the mapping data must come from the node: the shape
leaf paints it, internal nodes inherit the dominant child, and finalize
writes it into the vertex.  A mapping bug therefore shows up as a texture
scale collapse, not as a geometry error.

\subsection{The frame graph}

The main frame is roughly:
\begin{enumerate}
\item frame fences and pending command-buffer processing;
\item uploads and deferred destroys;
\item GPU cull for solid, water, SDF, bbox, vegetation;
\item depth prepass;
\item shadow cascades;
\item solid color pass;
\item water, vegetation, video, SDF passes;
\item post-process composition;
\item ImGui and present.
\end{enumerate}
The topology pass is not in this list: it runs on worker threads when a
chunk is built, not per frame.  A frame only consumes published spans.
